\chapter{后端库}\label{backend}

\section{\type {pdf}}

\topicindex{backend}
\topicindex{\PDF}

此库包含与 \PDF\ 后端相关的变量和函数。有关设置器的预期值的更多详细信息，
请参见 \ref{backendprimitives}。

\subsection{\type {mapfile}}

\topicindex{map files}
\topicindex{files+map}

\begin{typing}
pdf.mapfile(<string> map file)
pdf.mapline(<string> map line)
\end{typing}

这两个函数可用于替换从 \PDFTEX 继承的原语 \orm {pdfmapfile} 和 \orm {pdfmapline}。
它们接受一个字符串作为唯一参数，没有返回值。映射行中的第一个字符可以是
\type {-}、\type {+} 或 \type {=}，分别表示删除、添加或替换该行。
它们不是状态设置器，而是立即生效。

\subsection{{\type {[set|get}}}]

\topicindex{\PDF+trailer}
\topicindex{\PDF+catalog}
\topicindex{\PDF+info}

\libindex{setcatalog} \libindex{getcatalog}
\libindex{setinfo}    \libindex{getinfo}
\libindex{setnames}   \libindex{getnames}
\libindex{settrailer} \libindex{gettrailer}

这些函数补充了处理元数据的相应 \PDF\ 后端令牌列表。值类型为字符串，
它们在 \TEX\ 端设置的令牌寄存器写入之后直接写入 \PDF\ 文件。

\subsection{{\type {[set|get}}}]

\libindex{setpageattributes}  \libindex{getpageattributes}
\libindex{setpageresources}   \libindex{getpageresources}
\libindex{setpagesattributes} \libindex{getpagesattributes}

\topicindex{\PDF+page attributes}
\topicindex{\PDF+page resources}

这些函数补充了处理页面资源的相应 \PDF\ 后端令牌列表。这些变量与相应的
\PDF\ 后端令牌寄存器没有交互。它们在 \TEX\ 端设置的令牌寄存器写入之后
直接写入 \PDF\ 文件。

\subsection{{\type {[set|get}}}]

\libindex{setxformattributes} \libindex{getxformattributes}
\libindex{setxformresources}  \libindex{getxformresources}

\topicindex{\PDF+xform attributes}
\topicindex{\PDF+xform resources}

这些函数补充了处理可复用盒子和图像的相应 \PDF\ 后端令牌列表。这些变量
与相应的 \PDF\ 后端令牌寄存器没有交互。它们在 \TEX\ 端设置的令牌寄存器
写入之后直接写入 \PDF\ 文件。

\subsection{{\type {[set|get}version}}]

\topicindex{\PDF+version}

\libindex{getmajorversion} \libindex{setmajorversion}
\libindex{getminorversion} \libindex{setminorversion}

你可以设置输出的主版本号和次版本号。主版本号通常为~1，但当设置为~2时，
为了符合标准，某些数据将不会写入文件。你设置哪个次版本号取决于你使用的
\PDF\ 功能。这不在 \LUATEX\ 的控制范围内。

可以将主版本号设置为~2，但我们无法保证引擎能正确适配，因为据我们所知
没有公开免费的规范。此外，用户构造的注释不会被检查，而是直接传递给文件。
另一方面，引擎生成的 \PDF\ 相当简单，对版本的依赖不大。

\subsection{\type {getcreationdate}}

\topicindex{\PDF+date}

\libindex{getcreationdate}

此函数返回一个字符串，其中包含最终出现在 \PDF\ 文件中的日期格式，
本例中为：{\tttf [版本号]}。

\subsection{{\type {[set|get}inclusionerrorlevel} 和 \type {[set|get]ignoreunknownimages}}]

\topicindex{\PDF+options}

\libindex{getinclusionerrorlevel} \libindex{setinclusionerrorlevel}
\libindex{getignoreunknownimages} \libindex{setignoreunknownimages}

这些变量控制如何处理嵌入图像中的错误。它们是仿照 \PDFTEX\ 中的对应项
建模的。

\subsection{\type{[set|get]suppressoptionalinfo}, \type{[set|get]trailerid}, \type{[set|get]omitcidset}, \type{[set|get]omitinfo}, \type{[set|get]omitmediabox}, \type{[set|get]omitprocset}, \type{[set|get]ptexprefix}}

\topicindex{\PDF+options}
\topicindex{\PDF+trailer}

\libindex{getsuppressoptionalinfo} \libindex{setsuppressoptionalinfo}
\libindex{gettrailerid}            \libindex{settrailerid}
\libindex{getomitcidset}           \libindex{setomitcidset}
\libindex{getomitcharset}          \libindex{setomitcharset}
\libindex{getomitinfo}             \libindex{setomitinfo}
\libindex{getomitmediabox}         \libindex{setomitmediabox}
\libindex{getomitprocset}          \libindex{setomitprocset}
\libindex{getptexprefix}           \libindex{setptexprefix}

可选信息位集（一个数字）确定要刷新哪些类型的信息。默认情况下我们刷新所有。
更多详细信息请参见 \ref{sec:pdfextensions}。

你可以设置自己的 trailer id。这必须是一个包含有效 \PDF\ 数组内容和
校验和的字符串。

cidset、charset 和 info 标志（数字）会禁用所谓的 \type {CIDSet} 和
\type {CharSet} 条目的包含，这在针对众多 \PDF\ 子标准中的某些时可能很方便。
\type {ProcSet} 和 \type {PTEX} 前缀也是如此，设置此标志将强制使用
\type {_} 而不是 \type {.}。\footnote {在 info 字典中，句点在 2 版之前
是有效的，但在其他地方必须使用下划线。此前缀始于 \PDFTEX\ 的早期，
当时使用句点被认为是可以接受的。后来的规范对此进行了澄清。}

当省略它时，应通过页面属性选项提供 \type {MediaBox}，因为这是一个
必需字段。不进行任何检查。

\subsection{{\type {[set|get}compresslevel} 和 \type {[set|get]recompress}}]

\topicindex{\PDF+compression}

\libindex{getcompresslevel}    \libindex{setcompresslevel}
\libindex{getobjcompresslevel} \libindex{setobjcompresslevel}
\libindex{getrecompress}       \libindex{setrecompress}

这些函数设置流压缩级别。启用对象压缩时，多个对象将被打包到一个压缩流中，
从而节省空间。最小值为~0，最大值为~9。

当 recompression 设为~1时，压缩对象将被解压；当 compresslevel 大于零时，
它们将被重新压缩。这主要是一个调试功能，不应依赖。

\subsection{{\type {[set|get}gentounicode}}]

\topicindex{\PDF+unicode}

\libindex{getgentounicode} \libindex{setgentounicode}

此标志启用 tounicode 生成（如 \PDFTEX\ 中那样）。通常这些值由字体加载器提供。

\subsection{{\type {[set|get}decimaldigits}}]

\topicindex{\PDF+precision}

\libindex{getdecimaldigits} \libindex{setdecimaldigits}

这两个函数设置写入 \PDF 文件的浮点数精度。你可以设置任何值，但后端
不会低于~3或高于~6。

\subsection{{\type {[set|get}pkresolution}}]

\topicindex{\PDF+resolution}

\libindex{getpkresolution} \libindex{setpkresolution}

此设置器接受两个参数：分辨率和一个可选的零或一，指示这是否是固定的。
获取器返回这两个值。

\subsection{{\type {getlast[obj|link|annot}} 和 \type {getretval}}]

\topicindex{\PDF+objects}
\topicindex{\PDF+annotations}

\libindex{getlastobj}   \libindex{setlastobj}
\libindex{getlastlink}  \libindex{setlastlink}
\libindex{getlastannot} \libindex{setlastannot}
\libindex{getretval}

这些状态变量类似于传统上在 \TEX\ 端后端接口中使用的变量。

\subsection{\type {getmaxobjnum}}

\libindex{getmaxobjnum}
\libindex{getobjtype}
\libindex{getfontname}
\libindex{getfontobjnum}
\libindex{getfontsize}
\libindex{getxformname}

这些内省辅助函数主要在你自己构造 \PDF\ 对象时使用，例如需要有关
（将要）嵌入的字体的信息时。

\subsection{{\type {[set|get}origin}}]

\topicindex{\PDF+positioning}

\libindex{setorigin} \libindex{getorigin}

此函数用于设置水平和|/|或垂直偏移，这是一个传统的后端属性。

\begin{typing}
pdf.setorigin() -- sets both to 0pt
pdf.setorigin(tex.sp("1in")) -- sets both to 1in
pdf.setorigin(tex.sp("1in"),tex.sp("1in"))
\end{typing}

此函数的对应获取器返回两个值。

\subsection{{\type {[set|get}imageresolution}}]

\topicindex{\PDF+resolution}

\libindex{setimageresolution} \libindex{getimageresolution}

这两个函数与 imageresolution 相关，当图像本身未提供非零 x 或 y 分辨率时使用。

\subsection{{\type {[set|get}margin}}]

\topicindex{\PDF+margins}

\libindex{getlinkmargin}   \libindex{setlinkmargin}
\libindex{getdestmargin}   \libindex{setdestmargin}
\libindex{getthreadmargin} \libindex{setthreadmargin}
\libindex{getxformmargin}  \libindex{setxformmargin}
\libindex{getmarginmargin} \libindex{setmarginmargin}

这些函数可用于设置和检索添加到相应对象自然边界框的边距。

\subsection{{\type {get[pos|hpos|vpos}}}]

\topicindex{\PDF+positions}

\libindex{getpos}
\libindex{gethpos}
\libindex{getvpos}

这些函数获取输出页面上的当前位置，从其左下角开始测量。
返回值以缩放点为单位。

\begin{typing}
local h, v = pdf.getpos()
\end{typing}

\subsection{{\type {[has|get}matrix}}]

\topicindex{\PDF+matrix}

\libindex{getmatrix} \libindex{hasmatrix}

当前矩阵变换可通过 \type {getmatrix} 命令获得，该命令返回 6 个值：
\type {sx}、\type {rx}、\type {ry}、\type {sy}、\type {tx} 和 \type {ty}。
当应用了矩阵时，\type {hasmatrix} 函数返回 \type {true}。

\begin{typing}
if pdf.hasmatrix() then
    local sx, rx, ry, sy, tx, ty = pdf.getmatrix()
    -- do something useful or not
end
\end{typing}

\subsection{\type {print}}

\topicindex{\PDF+print to}

\libindex{print}

你可以从 \lpr {latelua} 调用内向 \PDF\ 文档打印一个字符串。
此函数不应在 \prm {directlua} 内部使用，除非你 {\it 确切地}知道自己在做什么。

\begin{typing}
pdf.print(<string> s)
pdf.print(<string> type, <string> s)
\end{typing}

可选参数可用于模拟 \PDF\ 字面量的行为：\type {type} 为 \type {direct}
或 \type {page}。

\subsection{\type {immediateobj}}

\topicindex{\PDF+objects}

\libindex{immediateobj}

此函数创建一个 \PDF\ 对象并立即将其写入 \PDF\ 文件。它是仿照
\PDFTEX 的 \prm {immediate} \orm {pdfobj} 原语建模的。所有函数变体都返回
新生成对象的对象编号。

\begin{typing}
<number> n =
    pdf.immediateobj(<string> objtext)
<number> n =
    pdf.immediateobj("file", <string> filename)
<number> n =
    pdf.immediateobj("stream", <string> streamtext, <string> attrtext)
<number> n =
    pdf.immediateobj("streamfile", <string> filename, <string> attrtext)
\end{typing}

第一个版本将 \type {objtext} 原样放入一个对象中。只有对象包装器是自动
生成的，但任何内部结构（如 \type{<< >>} 字典标记）都需要由用户提供。
第二个版本以关键字 \type {file} 作为第一个参数，将名为 \type {filename}
的文件内容原样放入对象中。第三个版本以关键字 \type {stream} 创建一个流
对象，并将 \type {streamtext} 原样放入流中。流长度自动计算。可选的
\type {attrtext} 放入该对象的字典中。第四个版本以关键字 \type {streamfile}
与第三个版本作用相同，只是从文件中原样读取流数据。

可以给出一个可选的第一个参数，使函数使用先前保留的 \PDF\ 对象。

\begin{typing}
<number> n =
    pdf.immediateobj(<integer> n, <string> objtext)
<number> n =
    pdf.immediateobj(<integer> n, "file", <string> filename)
<number> n =
    pdf.immediateobj(<integer> n, "stream", <string> streamtext, <string> attrtext)
<number> n =
    pdf.immediateobj(<integer> n, "streamfile", <string> filename, <string> attrtext)
\end{typing}

\subsection{\type{obj}}

\topicindex{\PDF+objects}

\libindex{obj}

此函数创建一个 \PDF\ 对象，该对象仅在被引用时（例如通过 \type {refobj()}）
才写入 \PDF\ 文件。

所有函数变体都返回新生成对象的对象编号，并且有两种独立的调用模式。
第一种模式是仿照 \PDFTEX 的 \orm {pdfobj} 原语建模的。

\begin{typing}
<number> n =
    pdf.obj(<string> objtext)
<number> n =
    pdf.obj("file", <string> filename)
<number> n =
    pdf.obj("stream", <string> streamtext, <string> attrtext)
<number> n =
    pdf.obj("streamfile", <string> filename, <string> attrtext)
\end{typing}

可以给出一个可选的第一个参数，使函数使用先前保留的 \PDF\ 对象。

\begin{typing}
<number> n =
    pdf.obj(<integer> n, <string> objtext)
<number> n =
    pdf.obj(<integer> n, "file", <string> filename)
<number> n =
    pdf.obj(<integer> n, "stream", <string> streamtext, <string> attrtext)
<number> n =
    pdf.obj(<integer> n, "streamfile", <string> filename, <string> attrtext)
\end{typing}

第二种模式接受一个带有键值对的参数表。

\begin{typing}
<number> n = pdf.obj {
    type           = <string>,
    immediate      = <boolean>,
    objnum         = <number>,
    attr           = <string>,
    compresslevel  = <number>,
    objcompression = <boolean>,
    file           = <string>,
    string         = <string>,
    nolength       = <boolean>,
}
\end{typing}

\type {type} 字段可以取 \type {raw} 和 \type {stream} 值，此字段为必填，
其他字段为可选（在约束范围内）。设置 \type {nolength} 时，不会向字典中
添加 \type {/Length} 条目。

注意：这种模式使 \type{obj} 看起来比实际更灵活：独立参数版本的约束仍然
适用，例如你不能同时使用 \type {string} 和 \type {file}。

\subsection{\type {refobj}}

\topicindex{\PDF+objects}

\libindex{refobj}

此函数是 \orm {pdfrefobj} 原语的 \LUA\ 版本，通过对象编号引用一个对象，
以便将该对象写入 \PDF\ 文件。

\begin{typing}
pdf.refobj(<integer> n)
\end{typing}

此函数在 \prm {directlua} 和 \lpr {latelua} 环境中都有效。
在 \prm {directlua} 内部，会创建一个新的 whatsit 节点 ``pdf_refobj，
该节点将在页面输出期间被标记为待刷新，然后该对象将在页面之后直接写入，
此时资源对象也会被写入 \PDF\ 文件。在 \lpr {latelua} 内部，该对象将被
标记为待刷新。

此函数没有返回值。

\subsection{\type {reserveobj}}

\topicindex{\PDF+objects}

\libindex{reserveobj}

此函数创建一个空的 \PDF\ 对象并返回其编号。

\begin{typing}
<number> n = pdf.reserveobj()
<number> n = pdf.reserveobj("annot")
\end{typing}

\subsection{\type {getpageref}}

\topicindex{\PDF+pages}

\libindex{getpageref}

可以通过此函数获取页面的对象编号。这可能是一个前向引用，因此当你请求
未来的页面时，确实会得到一个编号。

\begin{typing}
<number> n = pdf.getpageref(123)
\end{typing}

\subsection{\type {registerannot}}

\topicindex{\PDF+annotations}

\libindex{registerannot}

此函数将一个对象编号添加到当前页面的 \type {/Annots} 数组中，不做其他
任何事情。此函数只能在 \lpr {latelua} 内部使用。

\begin{typing}
pdf.registerannot(<number> objnum)
\end{typing}

\subsection{\type {newcolorstack}}

\topicindex{\PDF+color stack}

\libindex{newcolorstack}

此函数分配一个新的颜色堆栈并返回其 id。参数与类似的后端扩展原语相同。

\begin{typing}
pdf.newcolorstack("0 g","page",true) -- page|direct|origin
\end{typing}

\subsection{\type {setfontattributes}}

\topicindex{\PDF+fonts}

\libindex{setfontattributes}

此函数将强制向字体资源中添加一些额外代码。例如，它可用于向位图文件
添加自定义 \type {ToUnicode} 向量。

\begin{typing}
pdf.setfontattributes(<number> font id, <string> pdf code)
\end{typing}

\section{\type {pdfe}}

\subsection{简介}

\topicindex{\PDF+objects}

\topicindex{\PDF+analyze}
\topicindex{\PDF+\type{pdfe}}

\type {pdfe} 库取代了 \type {epdf} 库，提供了对 \PDF\ 文件的接口。
它使用与 \PDF\ 图像嵌入相同的代码。由 Paweł Jackowski 编写的
\type {pplib} 库取代了 \type {poppler}（派生自 \type {xpdf}）库。

\PDF\ 文件基本上是一棵对象树，通过字典（键/值）和数组（索引/值）
进入树中。有几个从根开始的最顶层字典可以更直接地访问。

尽管 \PDF\ 中的一切基本上都是对象，但我们只将其中一些封装为所谓的
userdata \LUA\ 对象。

\begin{longtable}{ll}
\toprule
\PDF          & \LUA \\
\midrule
null          & nil \\
boolean       & boolean \\
integer       & integer \\
float         & number \\
name          & string \\
string        & string \\
array         & array userdatum \\
dictionary    & dictionary userdatum \\
stream        & stream userdatum（带有相关字典） \\
reference     & reference userdatum \\
\bottomrule
\end{longtable}

常规获取器返回这些 \LUA\ 数据类型，但也可以获得更详细的信息。

\subsection{\type {open}}

\libindex {open}
\libindex {new}
\libindex {new}
\libindex {getstatus}
\libindex {close}
\libindex {unencrypt}

文档从文件或字符串加载

\begin{typing}
<pdfe document> = pdfe.open(filename)
<pdfe document> = pdfe.new(somestring,somelength)
\end{typing}

这样的文档用以下方式关闭：

\begin{typing}
pdfe.close(<pdfe document>)
\end{typing}

你可以通过以下方式检查文档是否成功打开：

\begin{typing}
pdfe.getstatus(<pdfe document>)
\end{typing}

返回的代码为：

\begin{longtable}{c l}
\toprule 值      & 说明 \\
\midrule
\type {-2} & 文档（仍然）受保护 \\
\type {-1} & 文档打开失败 \\
\type  {0} & 文档未加密 \\
\type  {1} & 文档已被解密 \\
\bottomrule
\end{longtable}

加密文档可以通过以下命令解密，其中两个密码参数都可以给出 \type {nil}：

\begin{typing}
pdfe.unencrypt(<pdfe document>,userpassword,ownerpassword)
\end{typing}

\subsection{\type {getsize}}

\libindex {getsize}
\libindex {getversion}
\libindex {getnofobjects}
\libindex {getnofpages}
\libindex {getmemoryusage}

成功打开的文档可以提供一些信息：

\begin{typing}
bytes = getsize(<pdfe document>
major, minor = getversion(<pdfe document>
n = getnofobjects(<pdfe document>
n = getnofpages(<pdfe document>
bytes, waste = getmemoryusage(<pdfe document>
\end{typing}

\subsection{{\type {get[catalog|trailer|info}}}]

\libindex {getcatalog}
\libindex {gettrailer}
\libindex {getinfo}

要访问文档结构，你从所谓的 catalog（目录）开始，它是一个字典：

\begin{typing}
<pdfe dictionary> = pdfe.getcatalog(<pdfe document>)
\end{typing}

另外两个根字典通过以下方式访问：

\begin{typing}
<pdfe dictionary> = pdfe.gettrailer(<pdfe document>)
<pdfe dictionary> = pdfe.getinfo(<pdfe document>)
\end{typing}

\subsection{\type {getpage}}

\libindex {getpage}
\libindex {getbox}

可以通过以下命令方便地访问特定页面，该命令返回一个字典。
第一个参数应是一个页面字典。

\begin{typing}
<pdfe dictionary> = pdfe.getpage(<pdfe document>,pagenumber)
\end{typing}

另一个便捷命令为你提供（通常是页面的）（边界）框，该框可以从文档本身
继承。有效框名的一个例子是 \type {MediaBox}。

\begin{typing}
box_dimens <array> = pdfe.getbox(<pdfe dictionary>,boxname)
\end{typing}

\subsection{{\type {get[string|integer|number|boolean|name}}, \type{type}}]

\libindex {getstring}
\libindex {getinteger}
\libindex {getnumber}
\libindex {getboolean}
\libindex {getname}
\libindex {type}

字典和数组中的常见值是字符串、整数、浮点数、布尔值和名称（它们也是
字符串），这些也是普通的 \LUA\ 对象：

\begin{typing}
s = getstring <pdfe array|dictionary>,index|key)
i = getinteger(<pdfe array|dictionary>,index|key)
n = getnumber <pdfe array|dictionary>,index|key)
b = getboolean(<pdfe array|dictionary>,index|key)
n = getname   <pdfe array|dictionary>,index|key)
s = type      <pdfe array|dictionary|document|reference|stream>
\end{typing}

\type {type} 返回一个描述对象类型的字符串，例如 "pdfe.array"、
"pdfe.dictionary"、"pdfe"、"pdfe.reference"、"pdfe.stream"。

\type {getstring} 函数有两个额外的变体：

\begin{typing}
s, h = getstring <pdfe array|dictionary>,index|key,false)
s    = getstring <pdfe array|dictionary>,index|key,true)
\end{typing}

第一次调用返回原始字符串加上一个布尔值，指示该字符串是否为十六进制编码。
第二次调用返回未编码的字符串。

\subsection{{\type {get[from}}}]

\libindex {getdictionary} \libindex {getfromdictionary}
\libindex {getarray}      \libindex {getfromarray}
\libindex {getstream}     \libindex {getfromstream}

通常你会在数组中使用索引，在字典中使用键，但字典也接受索引。
数组或字典的大小可以用通常的 \type {#} 运算符获得。

\begin{typing}
<pdfe dictionary>   = getdictionary(<pdfe array|dictionary>,index|key)
<pdfe array>        = getarray     <pdfe array|dictionary>,index|key)
<pdfe stream>,
<pdfe dictionary>   = getstream    <pdfe array|dictionary>,index|key)
\end{typing}

这些命令返回字典、数组和流，流是附加了一块数据的字典。

在介绍另一种访问模式之前，我们提到这些对象也以不同的方式提供访问，
例如以下是有效的：

\begin{typing}
print(pdfe.open("foo.pdf").Catalog.Type)
\end{typing}

在最顶层有 \type {Catalog}、\type {Info}、\type {Trailer} 和 \type {Pages}，
所以以下也是可以的：

\begin{typing}
print(pdfe.open("foo.pdf").Pages[1])
\end{typing}

\subsection{{\type {[open|close|readfrom|whole|}stream}}]

\libindex {openstream}
\libindex {closestream}
\libindex {readfromstream}
\libindex {readfromwholestream}

流有些特殊。当你的索引或键命中一个流时，你会得到一个流对象和字典对象。
字典你可以用通常的方式访问，而流有以下方法：

\begin{typing}
okay   = openstream(<pdfe stream>,[decode])
         closestream(<pdfe stream>
str, n = readfromstream(<pdfe stream>
str, n = readwholestream(<pdfe stream>,[decode])
\end{typing}

你可以分块读取，也可以一次性读取全部。分块读取时，你需要自己打开和关闭流。
\type {n} 值指示读取的长度。\type {decode} 参数控制是否对流数据进行解压。

与字典一样，你也可以用通常的 \LUA\ 方式访问流字典中的字段。当你
``调用 流时，你会得到内容。你可以传递一个布尔值，指示是否需要
解压流。

% pdfe.objectcodes      = objectcodes
% pdfe.stringcodes      = stringcodes
% pdfe.encryptioncodes  = encryptioncodes

\subsection{{\type {getfrom[dictionary|array}}}]

\libindex {getfromdictionary}
\libindex {getfromarray}

除了前面描述的接口之外，还有一个更低级别的接口可用。

\begin{typing}
key, type, value, detail = getfromdictionary(<pdfe dictionary>,index)
type, value, detail = getfromarray(<pdfe array>,index)
\end{typing}

\begin{longtable}{c l l l}
\toprule 类型       & 含义    & 值            & 详情 \\
\type {0}  & 无       & nil              & \\
\type {1}  & 空值     & nil              & \\
\type {2}  & 布尔值   & 布尔值           & \\
\type {3}  & 整数     & 整数             & \\
\type {4}  & 数字     & 浮点数           & \\
\type {5}  & 名称     & 字符串           & \\
\type {6}  & 字符串   & 字符串           & 十六进制 \\
\type {7}  & 数组     & 数组对象         & 大小 \\
\type {8}  & 字典     & 字典对象         & 大小 \\
\type {9}  & 流       & 流对象           & 字典大小 \\
\type {10} & 引用     & 整数             & \\
\bottomrule
\end{longtable}

\type {hex} 字符串（在 \PDF\ 文件中）由 \type {<>} 包围，
而普通字符串由括号界定。

\subsection{{\type {[dictionary|array}totable}}]

\libindex {dictionarytotable}
\libindex {arraytotable}

字典或表中的所有条目都可以通过以下命令获取，返回值是一个哈希表或索引表。

\begin{typing}
hash = dictionarytotable(<pdfe dictionary>
list = arraytotable(<pdfe array>
\end{typing}

你可以通过以下方式获取页面列表：

\begin{typing}
{ { <pdfe dictionary>, size, objnum }, ... } = pagestotable(<pdfe document>
\end{typing}

\subsection{\type {getfromreference}}

\libindex {getfromreference}

因为可能存在未解析的引用，可以通过以下方式解析引用对象：

\begin{typing}
type, <pdfe dictionary|array|stream>, detail = getfromreference(<pdfe reference>
\end{typing}

因此，作为第二个值，你会得到一个新的 \type {pdfe} userdata 对象，
你可以查询它。

\section{内存流}

\topicindex{\PDF+memory streams}

\libindex {new}

\type {pdfe.new} 接受三个参数：

\begin{longtable}{lp{8cm}}
\toprule 值           & 说明      \\
\midrule
\type {stream}  & 这是一个（用低级 \LUA\ 术语来说）轻量级 userdata
                        对象，即指向字节序列的指针 \\
\type {length}  & 这是流的长度，以字节为单位（流可以包含嵌入的
                        零字节） \\
\type {name}    & 可选，这是一个唯一标识符，用于对流进行哈希处理，
                        以便多个流不会占用更多内存 \\
\bottomrule
\end{longtable}

第三个参数是可选的。如果未给出，函数将像处理普通文件一样返回一个
\type {pdfe} 文档对象；否则它将返回一个文件名，可在其他地方（例如在
图像库中）用作伪文件来引用该流。

你也可以传递一个 \LUA\ 字符串，而不是轻量级 userdata 流（这实际上
不太稳定，但当你来自一个库时很方便），在这种情况下，给定的长度
（至多）是字符串长度。

该函数返回一个 \type {pdfe} 对象和一个字符串。该字符串可以在
\type {img} 库中代替文件名使用。当你将该对象用作图像时，需要防止它被
垃圾回收（例如通过将其存储在某处）。

内存流及其在图像库中的使用都是实验性的，可能会发生变化。如果你想知道
这在哪里可以使用：当你使用 \type {graphicmagick} 的 swiglib 库时，它可以
返回这样一个 userdata 对象。这允许在内存中进行转换并将结果直接传递给后端。
这在一遍式工作流中可能会节省一些运行时间。此功能目前不用于生产环境，
我们可能会提供更好的实现。

\section{\type {pdfscanner}}

\topicindex{\PDF+scanner}

\libindex {scan}

\type {pdfscanner} 库允许解释 \PDF\ 内容流和 \type {/ToUnicode}（cmap）流。
你可以从 \type {pdfe} 库获取这些流，如前面的章节所述。此库中只有一个
顶层函数：

\begin{typing}
pdfscanner.scan <pdfe stream>, <table> operatortable, <table> info)
pdfscanner.scan <pdfe array>, <table> operatortable, <table> info)
pdfscanner.scan <string>, <table> operatortable, <table> info)
\end{typing}

第一个参数应该是一个 \LUA\ 字符串，或者来自 \type {pdfe} 库的流或数组对象。
第二个参数 \type {operatortable} 应该是一个 \LUA\ 表，其中键是 \PDF\ 运算符
名称字符串，值是（由你定义的）用于处理这些运算符的 \LUA\ 函数。每当扫描器
在内容流中找到这些 \PDF\ 运算符之一时，就会调用这些函数。这些函数以两个
参数被调用：\type {scanner} 对象本身，以及作为 \type {pdfscanner.scan}
第三个参数传递的 \type {info} 表。

在内部，\type {pdfscanner.scan} 遍历流中的 \PDF\ 运算符，在内部栈上收集
操作数，直到找到一个 \PDF\ 运算符。如果该 \PDF\ 运算符的名称存在于
\type {operatortable} 中，则执行关联的函数。函数运行后（或没有函数可执行时），
内部操作数栈会被清空，为下一个运算符做准备，然后继续处理。

处理函数的 \type {scanner} 参数是必需的，因为它提供了各种方法来从内部
操作数栈中获取实际操作数。
操作数栈中获取实际操作数。

处理 \PDF\ 文档流的一个简单示例如下：

\begin{typing}
local operatortable = { }

operatortable.Do = function(scanner,info)
    local resources = info.resources
    if resources then
        local val     = scanner:pop()
        local name    = val[2]
        local xobject = resources.XObject
        print(info.space .. "Uses XObject " .. name)
        local resources = xobject.Resources
        if resources then
            local newinfo =  {
                space     = info.space .. "  ",
                resources = resources,
            }
            pdfscanner.scan(entry, operatortable, newinfo)
        end
    end
end

local function Analyze(filename)
    local doc = pdfe.open(filename)
    if doc then
        local pages = doc.Pages
        for i=1,#pages do
            local page = pages[i]
            local info = {
              space     = "  " ,
              resources = page.Resources,
            }
            print("Page " .. i)
         -- pdfscanner.scan(page.Contents,operatortable,info)
            pdfscanner.scan(page.Contents(),operatortable,info)
        end
    end
end

Analyze("foo.pdf")
\end{typing}

此示例遍历 \PDF 中的所有实际内容，并打印找到的 \type {XObject} 名称。
虽然这段代码演示了不少 \type {pdfe} 函数，但我们还是关注
\type {pdfscanner} 类型特有的代码。

自底向上看，以下行运行扫描器，第一个参数给出了 \PDF\ 页面的顶层内容。

第三个参数 \type {info} 包含两个条目：\type {space} 用于缩进打印输出，
\type {resources} 是必需的，以便嵌入的 \type {XForms} 能够找到自己的内容。

第二个参数 \type {operatortable} 为单个 \PDF\ 运算符 \type {Do} 定义了
一个处理函数。
函数 \type {Do} 打印当前 \type {XObject} 的名称，然后在该 \type {XObject}
实际上是一个 \type {/Form} 的条件下，为该对象的内容流启动一个新的扫描器。
那个嵌套扫描器使用新的 \type {info} 参数调用，其中 \type {space} 值已更新，
以便输出的缩进良好地嵌套；还带有一个新的 \type {resources} 字段，以帮助
下一次迭代正确处理任何其他嵌入的 \type {XObject}。

当然，这在实践中不是一个非常有用的示例，但对于演示 \type {pdfscanner}
来说，它的长度刚好。它只使用了一个 \type {scanner} 方法：
\type {scanner:pop()}。该函数弹出内部栈的顶部操作数，并返回一个 \LUA\ 表，
其中索引 1 处的对象是一个表示操作数类型的字符串，索引 2 处的对象是其值。

可能的操作数类型及关联的 \LUA\ 值类型列表如下：

\begin{longtable}{l l}
\toprule 类型           & \LUA 类型      \\
\midrule
\type{integer}  & <number>  \\
\type{real}     & <number>  \\
\type{boolean}  & <boolean> \\
\type{name}     & <string>  \\
\type{operator} & <string>  \\
\type{string}   & <string>  \\
\type{array}    & <table>   \\
\type{dict}     & <table>   \\
\bottomrule
\end{longtable}

对于 \type {integer} 或 \type {real}，值始终是一个 \LUA{}（浮点）数字。
对于 \type {name}，前导斜杠总是被去掉。

对于 \type {string}，请注意 \PDF\ 实际上在 \PDF\ 文档的不同部分支持
不同类型的字符串（具有不同的编码），因此你可能需要对某些结果重新编码；
\type {pdfscanner} 始终输出字节流而不进行任何重新编码。
\type {pdfscanner} 不区分字面字符串和十六进制字符串（十六进制值会被解码），
并且它将内嵌图像的流数据视为一个字符串，即 \type {EI} 的唯一操作数。

对于 \type {array}，表内容是 \type {pop} 返回值的列表；对于 \type {dict}，
表键是 \PDF\ 名称字符串，值是 \type {pop} 返回值。

\libindex{pop}
\libindex{popnumber}
\libindex{popname}
\libindex{popstring}
\libindex{poparray}
\libindex{popdictionary}
\libindex{popboolean}
\libindex{done}

还有一些其他方法可以供你向 \type {scanner} 调用：

\begin{longtable}{l p{5cm}}
\toprule 方法               & 说明 \\
\midrule
\type{pop}           & 见上文 \\
\type{popnumber}     & 仅返回 \type {real} 或 \type {integer} 的值 \\
\type{popname}       & 仅返回 \type {name} 的值 \\
\type{popstring}     & 仅返回 \type {string} 的值 \\
\type{poparray}      & 仅返回 \type {array} 的值 \\
\type{popdictionary} & 仅返回 \type {dict} 的值 \\
\type{popboolean}    & 仅返回 \type {boolean} 的值 \\
\type{done}          & 中止此 \type {scan()} 调用的进一步处理 \\
\bottomrule
\end{longtable}

\type {pop*} 是便捷函数，当你事先知道操作数的类型时（在 \PDF 中你通常
确实知道）会很方便。例如，\type {Do} 函数本可以使用
\type {local name = scanner:popname()}，因为 \type {Do} 运算符的唯一
操作数始终是一个 \PDF\ 名称对象。

\type {done} 函数允许你在已经了解了所有想了解的内容后中止流的处理。
这在解析 \type {/ToUnicode} 时很方便，因为通常有一些你不感兴趣的尾部
垃圾数据。如果没有 \type {done}，处理只会在流的末尾结束，可能会浪费
\CPU\ 周期。

{\em 我们保留了旧的名称 \type {popNumber}、\type {popName}、
\type {popString}、\type {popArray}、\type {popDict} 和 \type {popBool}。}
